\uuid{4tgT}
\exo7id{3958}
\titre{exo7 3958}
\auteur{quercia}
\organisation{exo7}
\datecreate{2010-03-11}
\isIndication{false}
\isCorrection{false}
\chapitre{Dérivabilité des fonctions réelles}
\sousChapitre{Autre}
\cpge{mpsi-mp2i.derivabilite.c, pcsi.derivabilite.c, ptsi.derivabilite.c, mpsi-mp2i.procedes-sommatoires-discrets.a, pcsi.series-numeriques.a, ptsi.series-numeriques.a, bcpst1.derivation.b, mp-mpi.series-numeriques-vectorielles.b, pc.series-numeriques, psi.series-numeriques.a, pt.series-numeriques.a, bcpst2.series-reelles}
\module{Analyse}
\niveau{L1}
\difficulte{}

\contenu{
\question{
Pour $k\in\N$, $k\ge 2$, appliquer le théorème des accroissements finis à
$x \mapsto\ln(\ln x)$ sur $[k,k+1]$.
En déduire que la série de terme général $\frac1{k\ln k}$ est divergente.
}
}
